% A4 is 21 x 29.7 cm
% leave 2.5 cm margins

\documentclass[a4paper,12pt]{article}%
\usepackage[dvips]{graphicx}
\usepackage{amsmath}
\usepackage{amsfonts}
\usepackage{amssymb}%
\setcounter{MaxMatrixCols}{30}
%TCIDATA{OutputFilter=latex2.dll}
%TCIDATA{Version=5.50.0.2953}
%TCIDATA{CSTFile=C:\swp55\styles\article\Joan Decker article.cst}
%TCIDATA{Created=Saturday, March 24, 2007 21:51:26}
%TCIDATA{LastRevised=Sunday, March 25, 2007 19:31:13}
%TCIDATA{<META NAME="GraphicsSave" CONTENT="32">}
%TCIDATA{<META NAME="SaveForMode" CONTENT="1">}
%TCIDATA{BibliographyScheme=Manual}
%TCIDATA{<META NAME="DocumentShell" CONTENT="Joan Decker\letter_jd_A4">}
%BeginMSIPreambleData
\providecommand{\U}[1]{\protect\rule{.1in}{.1in}}
%EndMSIPreambleData
\setlength{\textwidth}{16 true cm}
\setlength{\textheight}{24.7 true cm}
\setlength{\headheight}{0 true cm}
\setlength{\headsep}{0 true cm}
\setlength{\oddsidemargin}{0 true cm}
\setlength{\evensidemargin}{0 true cm}
\setlength{\topmargin}{0 true cm}
\begin{document}

\title{LUKE Simulations of LHCD in ITER}
\author{Joan Decker and Yves Peysson}
\date{\today}
\maketitle

\section{Introduction}

The LUKE ray-tracing/Fokker-Planck package is used to perform simulations of
LHCD in ITER. An up/down symmetric ITER Scenario 4 equilibrium is used. The
equilibrium data were provided by P. Bonoli for ITER simulations presented at
the IAEA conference in 2006. Figures representing the equilibrium are
available in Section 2. In section 3 a scan in $N_{\Vert}$ is performed for
the wave frequencies $f=5.0$ GHz and $f=3.7$ GHz, respectively. A\ poloidal
scan of ray lauching positions is presented in section 4. The results are
discussed in Section 5.

\section{Equilibrium}

The magnetic equilibrium is plotted in Fig. \ref{equil} and the profiles in
Fig. \ref{prof}.

\section{$N_{\Vert}$ scan at $f=5.0$ GHz and $f=3.7$ GHz}

For each scan, six rays starting from the outboard side are launched into the
plasma and propagated by the ray-tracing code C3PO. The six rays correspond to
the positive and negative lobes at the three different vertical launching
positions $Z=-0.5/0/0.5$ m. The total power in $P=30$ MW. The current
deposition profiles and total current corresponding to a scan with $N_{\Vert
}=1.9/2.0/2.1$ are shown in Fig. \ref{Jrf5p0} and \ref{Jrf3p7} corresponding
to the frequencies $f=5.0$ GHz and $f=3.7$ GHz, respectively. It is found that
the current drive efficiency increases for lower $N_{\Vert}$.\ This trend is
expected since the velocity $v_{\Vert}=c/N_{\Vert}$ of resonant electrons is
larger for lower $N_{\Vert}$. We note that the trend holds in all cases except
for $f=3.7$ GHz and $N_{\Vert}=1.9$ where the driven current is much lower
than expected.

For a given initial value of $N_{\Vert}$ (except the case $N_{\Vert}=1.9$),
the current drive efficiency is significantly higher for $f=3.7$ GHz than for
$f=5.0$ GHz. To interprete this results the rays are showed in Fig.
\ref{n2p0_rays} for $N_{\Vert}=2.0$ and each frequency. The value of
$N_{\Vert}$ at the location of maximum power deposition is displayed in Fig.
\ref{n2p0_Npar}. We observe that the ray propagation is different for the two
frequencies. When $f=3.7$ GHz, the rays experience more refraction and a
larger $N_{\Vert}$ downshift and thus propagate further into the plasma. This
explains the larger current drive efficiency.

In order to explain what happens for $f=3.7$ GHz and $N_{\Vert}=1.9$, the rays
are shown in Fig. \ref{n1p9_rays} for $N_{\Vert}=1.9$ and each frequency. The
value of $N_{\Vert}$ at the location of maximum power deposition is displayed
in Fig. \ref{n1p9_Npar}. We observe that for $f=3.7$ GHz two of the three rays
experience so much refraction that they propagate back towards the edge
without being damped. After reflection on the edge they experience a very
large $N_{\Vert}$ upshift leading to very inefficient current drive localized
close to the edge.

\section{Poloidal scan}

Rays were launched from 16 different, uniformly distributed, poloidal
locations with $f=5.0$ GHz and initial $N_{\Vert}=2.0$. These rays are shown
on Fig. \ref{raypol}. Rays are counted counted-clockwise starting with the ray
launched from the horizontal outboard midplane. On Fig. \ref{rhopol} and Fig.
\ref{nparpol}, the damping radial location and corresponding $N_{\Vert}$ value
are displayed, respectively, as a function of the ray number. It is found that
rays lauched from the horizontal midplane or slightly above (either on the LFS
or the HFS) experience a larger downshift and propagate further into the
plasma than rays lauched from the top or bottom of the plasma. Because these
rays at absorbed at a lower value of $N_{\Vert}$ they also have the potential
for more efficient current drive.

\section{Discussions and Conclusions}

Concerning the poloidal location of ray launching, it is demonstrated that
rays experience a larger $N_{\Vert}$ downshift and thus propagate further into
the plasma if they are launched on or slightly above the horizontal midplane.
There are no significant differences between LFS and HFS launching.

For a given frequency, rays launched with a lower $N_{\Vert}$ propagate
further into the plasma and lead to a larger current drive efficiency. This
trend holds when lowering $N_{\Vert}$ until the rays are refracted back to the
edge before being damped, in which case the contribution of such ray to
driving current drops considerably. For a given initial $N_{\Vert}$, rays
launched at the frequency $f=3.7$ GHz experience a stronger downshift and thus
drive more current than with $f=5.0$ GHz. This also means that the limit in
$N_{\Vert}$ for which the ray is reflected back to the edge before being
damped, which must be avoided, is higher in the case where $f=3.7$ GHz.

More generally, this limit in $N_{\Vert}$ results from the balance between
refraction - which depends upon the density profile - and damping - which
depends upon the temperature profile. The objective is to keep $N_{\Vert}$
high enough such that the ray is damped before being refracted, but low enough
to maintain a good current drive efficiency.

The maximum current drive efficiency that can be obtained is similar with
$f=3.7$ GHz and with $f=5.0$ GHz, but it occurs at a higher value of
$N_{\Vert}$ in the $f=3.7$ GHz case.%

%TCIMACRO{\FRAME{ftbhFU}{6.0494in}{4.5558in}{0pt}{\Qcb{Magnetic equilibrium of
%ITER scenario 4. The red and blue lines are contours of constant $\psi$ and
%$\theta$, respectively.}}{\Qlb{equil}}{iter_scen4_iaea06_equil.eps}%
%{\special{ language "Scientific Word";  type "GRAPHIC";
%maintain-aspect-ratio TRUE;  display "USEDEF";  valid_file "F";
%width 6.0494in;  height 4.5558in;  depth 0pt;  original-width 5.834in;
%original-height 4.3863in;  cropleft "0";  croptop "1";  cropright "1";
%cropbottom "0";
%filename 'EQUIL/figures/ITER_Scen4_IAEA06_equil.eps';file-properties "XNPEU";}%
%}}%
%BeginExpansion
\begin{figure}
[tbh]
\begin{center}
\includegraphics[
height=4.5558in,
width=6.0494in
]%
{EQUIL/figures/ITER_Scen4_IAEA06_equil.eps}%
\caption{Magnetic equilibrium of ITER scenario 4. The red and blue lines are
contours of constant $\psi$ and $\theta$, respectively.}%
\label{equil}%
\end{center}
\end{figure}
%EndExpansion%
%TCIMACRO{\FRAME{ftbhFU}{6.0494in}{3.6426in}{0pt}{\Qcb{Radial profiles of
%normalized poloidal magnetic flux, electron temperature, electron density, ion
%temperatures, ion densities, and effective charge for the ITER Scenario 4
%equilibrium.}}{\Qlb{prof}}{iter_scen4_iaea06_prof.eps}%
%{\special{ language "Scientific Word";  type "GRAPHIC";
%maintain-aspect-ratio TRUE;  display "USEDEF";  valid_file "F";
%width 6.0494in;  height 3.6426in;  depth 0pt;  original-width 9.9998in;
%original-height 6.0027in;  cropleft "0";  croptop "1";  cropright "1";
%cropbottom "0";
%filename 'EQUIL/figures/ITER_Scen4_IAEA06_prof.eps';file-properties "XNPEU";}%
%}}%
%BeginExpansion
\begin{figure}
[tbh]
\begin{center}
\includegraphics[
height=3.6426in,
width=6.0494in
]%
{EQUIL/figures/ITER_Scen4_IAEA06_prof.eps}%
\caption{Radial profiles of normalized poloidal magnetic flux, electron
temperature, electron density, ion temperatures, ion densities, and effective
charge for the ITER Scenario 4 equilibrium.}%
\label{prof}%
\end{center}
\end{figure}
%EndExpansion
%

%TCIMACRO{\FRAME{ftbhFU}{3.6002in}{2.3999in}{0pt}{\Qcb{Current density profile
%and total driven current for $f=5.0$ GHz and $N_{\Vert}=1.9$ (red), $N_{\Vert
%}=2.0$ (blue), and $N_{\Vert}=2.1$ (green). }}{\Qlb{Jrf5p0}}%
%{iter_scen4_iaea06_f5p0_npar_jrf.eps}{\special{ language "Scientific Word";
%type "GRAPHIC";  maintain-aspect-ratio TRUE;  display "USEDEF";
%valid_file "F";  width 3.6002in;  height 2.3999in;  depth 0pt;
%original-width 5.834in;  original-height 4.3863in;  cropleft "0";
%croptop "1";  cropright "1";  cropbottom "0";
%filename 'figures/ITER_Scen4_IAEA06_f5p0_Npar_JRF.eps';file-properties "XNPEU";}%
%}}%
%BeginExpansion
\begin{figure}
[tbh]
\begin{center}
\includegraphics[
height=2.3999in,
width=3.6002in
]%
{figures/ITER_Scen4_IAEA06_f5p0_Npar_JRF.eps}%
\caption{Current density profile and total driven current for $f=5.0$ GHz and
$N_{\Vert}=1.9$ (red), $N_{\Vert}=2.0$ (blue), and $N_{\Vert}=2.1$ (green). }%
\label{Jrf5p0}%
\end{center}
\end{figure}
%EndExpansion
%

%TCIMACRO{\FRAME{ftbhFU}{3.6002in}{2.3999in}{0pt}{\Qcb{Current density profile
%and total driven current for $f=3.7$ GHz and $N_{\Vert}=1.9$ (red), $N_{\Vert
%}=2.0$ (blue), and $N_{\Vert}=2.1$ (green). }}{\Qlb{Jrf3p7}}%
%{iter_scen4_iaea06_f3p7_npar_jrf.eps}{\special{ language "Scientific Word";
%type "GRAPHIC";  maintain-aspect-ratio TRUE;  display "USEDEF";
%valid_file "F";  width 3.6002in;  height 2.3999in;  depth 0pt;
%original-width 5.834in;  original-height 4.3863in;  cropleft "0";
%croptop "1";  cropright "1";  cropbottom "0";
%filename 'figures/ITER_Scen4_IAEA06_f3p7_Npar_JRF.eps';file-properties "XNPEU";}%
%}}%
%BeginExpansion
\begin{figure}
[tbh]
\begin{center}
\includegraphics[
height=2.3999in,
width=3.6002in
]%
{figures/ITER_Scen4_IAEA06_f3p7_Npar_JRF.eps}%
\caption{Current density profile and total driven current for $f=3.7$ GHz and
$N_{\Vert}=1.9$ (red), $N_{\Vert}=2.0$ (blue), and $N_{\Vert}=2.1$ (green). }%
\label{Jrf3p7}%
\end{center}
\end{figure}
%EndExpansion
%

%TCIMACRO{\FRAME{ftbhFU}{6.0494in}{4.5558in}{0pt}{\Qcb{Poloidal ray
%trajectories of the three rays corresponding to the positive lobe and three
%different poloidal launching positions. The rays are calculated for $N_{\Vert
%}=2.0$, and $f=5.0$ GHz (blue) or $f=3.7$ GHz (red).}}{\Qlb{n2p0_rays}%
%}{iter_scen4_iaea06_npar20_rays.eps}{\special{ language "Scientific Word";
%type "GRAPHIC";  maintain-aspect-ratio TRUE;  display "USEDEF";
%valid_file "F";  width 6.0494in;  height 4.5558in;  depth 0pt;
%original-width 5.834in;  original-height 4.3863in;  cropleft "0";
%croptop "1";  cropright "1";  cropbottom "0";
%filename 'WAVE/C3PO/figures/ITER_Scen4_IAEA06_Npar20_rays.eps';file-properties "XNPEU";}%
%}}%
%BeginExpansion
\begin{figure}
[tbh]
\begin{center}
\includegraphics[
height=4.5558in,
width=6.0494in
]%
{WAVE/C3PO/figures/ITER_Scen4_IAEA06_Npar20_rays.eps}%
\caption{Poloidal ray trajectories of the three rays corresponding to the
positive lobe and three different poloidal launching positions. The rays are
calculated for $N_{\Vert}=2.0$, and $f=5.0$ GHz (blue) or $f=3.7$ GHz (red).}%
\label{n2p0_rays}%
\end{center}
\end{figure}
%EndExpansion
%

%TCIMACRO{\FRAME{ftbhFU}{3.6002in}{2.3999in}{0pt}{\Qcb{Value of $N_{\Vert}$ at
%the location of maximum power deposition for the six rays shown in Fig.
%\ref{n2p0_rays}. Rays \#1,2,3 correspond to $f=5.0$ GHz and $Z=-0.5/0/0.5$ m,
%while rays \#4,5,6 correspond to the same vertical launching positions with
%$f=3.7$ GHz.}}{\Qlb{n2p0_Npar}}{iter_scen4_iaea06_npar20_npars.eps}%
%{\special{ language "Scientific Word";  type "GRAPHIC";
%maintain-aspect-ratio TRUE;  display "USEDEF";  valid_file "F";
%width 3.6002in;  height 2.3999in;  depth 0pt;  original-width 5.834in;
%original-height 4.3863in;  cropleft "0";  croptop "1";  cropright "1";
%cropbottom "0";
%filename 'WAVE/C3PO/figures/ITER_Scen4_IAEA06_Npar20_Npars.eps';file-properties "XNPEU";}%
%}}%
%BeginExpansion
\begin{figure}
[tbh]
\begin{center}
\includegraphics[
height=2.3999in,
width=3.6002in
]%
{WAVE/C3PO/figures/ITER_Scen4_IAEA06_Npar20_Npars.eps}%
\caption{Value of $N_{\Vert}$ at the location of maximum power deposition for
the six rays shown in Fig. \ref{n2p0_rays}. Rays \#1,2,3 correspond to $f=5.0$
GHz and $Z=-0.5/0/0.5$ m, while rays \#4,5,6 correspond to the same vertical
launching positions with $f=3.7$ GHz.}%
\label{n2p0_Npar}%
\end{center}
\end{figure}
%EndExpansion
%

%TCIMACRO{\FRAME{ftbhFU}{6.0494in}{4.5558in}{0pt}{\Qcb{Poloidal ray
%trajectories of the three rays corresponding to the positive lobe and three
%different poloidal launching positions. The rays are calculated for $N_{\Vert
%}=1.9$, and $f=5.0$ GHz (blue) or $f=3.7$ GHz (red).}}{\Qlb{n1p9_rays}%
%}{iter_scen4_iaea06_npar19_rays.eps}{\special{ language "Scientific Word";
%type "GRAPHIC";  maintain-aspect-ratio TRUE;  display "USEDEF";
%valid_file "F";  width 6.0494in;  height 4.5558in;  depth 0pt;
%original-width 5.834in;  original-height 4.3863in;  cropleft "0";
%croptop "1";  cropright "1";  cropbottom "0";
%filename 'WAVE/C3PO/figures/ITER_Scen4_IAEA06_Npar19_rays.eps';file-properties "XNPEU";}%
%}}%
%BeginExpansion
\begin{figure}
[tbh]
\begin{center}
\includegraphics[
height=4.5558in,
width=6.0494in
]%
{WAVE/C3PO/figures/ITER_Scen4_IAEA06_Npar19_rays.eps}%
\caption{Poloidal ray trajectories of the three rays corresponding to the
positive lobe and three different poloidal launching positions. The rays are
calculated for $N_{\Vert}=1.9$, and $f=5.0$ GHz (blue) or $f=3.7$ GHz (red).}%
\label{n1p9_rays}%
\end{center}
\end{figure}
%EndExpansion
%

%TCIMACRO{\FRAME{ftbhFU}{3.6002in}{2.3999in}{0pt}{\Qcb{Value of $N_{\Vert}$ at
%the location of maximum power deposition for the six rays shown in Fig.
%\ref{n1p9_rays}. Rays \#1,2,3 correspond to $f=5.0$ GHz and $Z=-0.5/0/0.5$ m,
%while rays \#4,5,6 correspond to the same vertical launching positions with
%$f=3.7$ GHz.}}{\Qlb{n1p9_Npar}}{iter_scen4_iaea06_npar19_npars.eps}%
%{\special{ language "Scientific Word";  type "GRAPHIC";
%maintain-aspect-ratio TRUE;  display "USEDEF";  valid_file "F";
%width 3.6002in;  height 2.3999in;  depth 0pt;  original-width 5.834in;
%original-height 4.3863in;  cropleft "0";  croptop "1";  cropright "1";
%cropbottom "0";
%filename 'WAVE/C3PO/figures/ITER_Scen4_IAEA06_Npar19_Npars.eps';file-properties "XNPEU";}%
%}}%
%BeginExpansion
\begin{figure}
[tbh]
\begin{center}
\includegraphics[
height=2.3999in,
width=3.6002in
]%
{WAVE/C3PO/figures/ITER_Scen4_IAEA06_Npar19_Npars.eps}%
\caption{Value of $N_{\Vert}$ at the location of maximum power deposition for
the six rays shown in Fig. \ref{n1p9_rays}. Rays \#1,2,3 correspond to $f=5.0$
GHz and $Z=-0.5/0/0.5$ m, while rays \#4,5,6 correspond to the same vertical
launching positions with $f=3.7$ GHz.}%
\label{n1p9_Npar}%
\end{center}
\end{figure}
%EndExpansion
%

%TCIMACRO{\FRAME{ftbhFU}{6.0494in}{4.5558in}{0pt}{\Qcb{Poloidal trajectories of
%LH rays. The black lines represent the ray portion where damping is
%significant.}}{\Qlb{raypol}}{iter_scen4_iaea06_poloidal_rays.eps}%
%{\special{ language "Scientific Word";  type "GRAPHIC";
%maintain-aspect-ratio TRUE;  display "USEDEF";  valid_file "F";
%width 6.0494in;  height 4.5558in;  depth 0pt;  original-width 5.834in;
%original-height 4.3863in;  cropleft "0";  croptop "1";  cropright "1";
%cropbottom "0";
%filename 'Poloidal/WAVE/figures/ITER_Scen4_IAEA06_poloidal_rays.eps';file-properties "XNPEU";}%
%}}%
%BeginExpansion
\begin{figure}
[tbh]
\begin{center}
\includegraphics[
height=4.5558in,
width=6.0494in
]%
{Poloidal/WAVE/figures/ITER_Scen4_IAEA06_poloidal_rays.eps}%
\caption{Poloidal trajectories of LH rays. The black lines represent the ray
portion where damping is significant.}%
\label{raypol}%
\end{center}
\end{figure}
%EndExpansion
%

%TCIMACRO{\FRAME{ftbhFU}{3.6002in}{2.3999in}{0pt}{\Qcb{Radial position of the
%peak location $\rho^{\text{abs}}$ in the power deposition profile.
%$\rho^{\text{abs}}\left(  \psi\right)  $ corresponds to $r/a$ where the flux
%surface $\psi$ crosses the outboard horizontal midplace.}}{\Qlb{rhopol}%
%}{iter_scen4_iaea06_poloidal_rhoabs.eps}%
%{\special{ language "Scientific Word";  type "GRAPHIC";
%maintain-aspect-ratio TRUE;  display "USEDEF";  valid_file "F";
%width 3.6002in;  height 2.3999in;  depth 0pt;  original-width 5.834in;
%original-height 4.3863in;  cropleft "0";  croptop "1";  cropright "1";
%cropbottom "0";
%filename 'Poloidal/WAVE/figures/ITER_Scen4_IAEA06_poloidal_rhoabs.eps';file-properties "XNPEU";}%
%}}%
%BeginExpansion
\begin{figure}
[tbh]
\begin{center}
\includegraphics[
height=2.3999in,
width=3.6002in
]%
{Poloidal/WAVE/figures/ITER_Scen4_IAEA06_poloidal_rhoabs.eps}%
\caption{Radial position of the peak location $\rho^{\text{abs}}$ in the power
deposition profile. $\rho^{\text{abs}}\left(  \psi\right)  $ corresponds to
$r/a$ where the flux surface $\psi$ crosses the outboard horizontal midplace.}%
\label{rhopol}%
\end{center}
\end{figure}
%EndExpansion
%

%TCIMACRO{\FRAME{ftbhFU}{3.6002in}{2.3999in}{0pt}{\Qcb{Value of $N_{\Vert}$
%along the ray at the position of maximum damping $\rho^{\text{abs}}$.}%
%}{\Qlb{nparpol}}{iter_scen4_iaea06_poloidal_nparabs.eps}%
%{\special{ language "Scientific Word";  type "GRAPHIC";
%maintain-aspect-ratio TRUE;  display "USEDEF";  valid_file "F";
%width 3.6002in;  height 2.3999in;  depth 0pt;  original-width 5.834in;
%original-height 4.3863in;  cropleft "0";  croptop "1";  cropright "1";
%cropbottom "0";
%filename 'Poloidal/WAVE/figures/ITER_Scen4_IAEA06_poloidal_Nparabs.eps';file-properties "XNPEU";}%
%}}%
%BeginExpansion
\begin{figure}
[tbh]
\begin{center}
\includegraphics[
height=2.3999in,
width=3.6002in
]%
{Poloidal/WAVE/figures/ITER_Scen4_IAEA06_poloidal_Nparabs.eps}%
\caption{Value of $N_{\Vert}$ along the ray at the position of maximum damping
$\rho^{\text{abs}}$.}%
\label{nparpol}%
\end{center}
\end{figure}
%EndExpansion



\end{document}